    \chapter{Comentarios sobre la revista}
    
\section{Objetivos}

La existencia de esta revista buscaría cumplir 4 objetivos:

\begin{enumerate}
    \item Familiarizar a los alumnos con el proceso de publicación de artículos.
    \item Dar a los alumnos material extra el cual presentar en sus exposiciones.
        \subitem Además de presentar un enlace al archivo digital en las presentaciones, podrían imprimirse algunas copias (3-4) para las presentaciones.
    \item Aumentar el entusiasmo de los estudiantes al permitirles ver su trabajo en un formato similar al de los formatos académicos.
    \item Generar una biblioteca de los Proyectos Integradores llevados a cabo por los alumnos, para así permitir su futura consulta por nuevas generaciones.
        \subitem Podría considerarse incluir tesinas realizadas por alumnos de sexto y décimo cuatrimestre en sus estadías, ya sea dentro del mismo archivo de la revista o como un apartado extra en el sitio web.
\end{enumerate}

\section{Diseño de la revista}

Hasta ahora, las inspiraciones para la revista han sido \href{https://www.pentagram.com/work/scientific-american?rel=sector&rel-id=3}{un rediseño de la revista Scientific American}, la revista \href{https://www.frontiersin.org/journals/sustainability/articles/10.3389/frsus.2026.1818134/pdf}{Frontiers}, \href{https://www.sciencedirect.com/science/article/pii/S1198743X26000650/pdfft?pid=1-s2.0-S1198743X26000650-main.pdf}{Elsevier} y algunas ideas en \href{https://pin.it/3s68QTUbZ}{Pinterest}.

Si bien aún pueden realizarse mejoras, se esperará a recibir retroalimentación de docentes y alumnos para aplicarlas. Asímismo, en un futuro también se aceptarán recomendaciones de alumnos y docentes para mejorar el diseño de la revista.

\section{Proceso de envio}

Para familiarizar a los alumnos con el proceso de publicación, se podría establecer un formato de envio el cual los alumnos deberán seguir para la publicación de sus artículos. El formato podría ser determinado en base a recomendaciones de docentes para que sea lo más similar posible al estándar académico.

A continuación se presentan algunos de los aspectos a definir en el formato, así como recomendaciones:

\begin{itemize}
    \item Tipo de archivo
        \subitem Google Docs o Word
    \item Características de la tipografía
        \subitem Fuente: Times New Roman o Arial
        \subitem Tamaño de fuente: 12pt
        \subitem Interlineado: 1.5
        \subitem Titulos: En negritas
        \subitem Nombres científicos en itálica
    \item Características de la hoja
        \subitem Márgenes sin restricciones
        \subitem Sin encabezado ni pie de página
    \item Figuras
        \subitem Enviar archivos de imagen aparte
        \subitem En lugar de adjuntar las figuras, específicar el lugar preferido donde se quiere que se muestren, para así aligerar el archivo de texto enviado
    \item Tablas
        \subitem De realizarse en hojas de cálculo (Excel, Sheets, etc.), adjuntar tales archivos
    \item Bibliografía
        \subitem Formato APA7
\end{itemize}

Dado que los artículos que se publicarán son producto de una materia con docente asignado, tal docente llevaría a cabo la revisión de los textos. Aún así, podría existir un proceso de revisión por pares realizado por alumnos de mayor grado. Tal ejercicio podría expandir los beneficios académicos y pedagógicos de la revista al acercar a más alumnos al mundo de la academia y sus procesos.

La plantilla de LaTeX mediante la cual se maqueta la revista permite preparar los número unicamente copiando y pegando lo escrito por los alumnos en lugares determinados, por lo que la preparación de la revista sería algo sencillo que podría realizarse por una persona. Personalmente disfruto del proceso, por lo que me gustaría ser el encargado de tal área, sin embargo, dado que LaTeX es una herramienta útil para el procesamiento de archivos, estaría abierto a ceder tal proceso a otros alumnos si estos se muestran interesados a aprender sobre el programa. 

\section{Tipos de artículos}

La revista tendría 3 tipos de artículos:

\begin{itemize}
    \item Artículo de Investigación
    \item Revisión/Review
    \item Especial\footnote{El nombre puede cambiar en un futuro}
\end{itemize}

Los primeros 2 son aquellos que los alumnos pueden llevar a cabo en la materia de Proyecto Integrador, por lo que solo sería necesario que lo especifiquen al momento de enviar su artículo.

El último es un tipo de artículo propuesto para expandir la participación de los alumnos en la revista. Consistiría en artículos donde su formato sería más libre, aceptando artículos de divulgación, ensayos, producciones artísticas (poemas e ilustraciones, por ejemplo), etc., y donde las temáticas sean más variadas, tocando áreas de la biotecnología menos exploradas en la universidad, tales como las ciencias sociales, por ejemplo.

Los artículos de esta área los podrían enviar alumnos de cualquier cuatrimestre de manera voluntaria. Ya que estos trabajos no formarían parte de alguna asignatura, su proceso de revisión se establecería dependiendo del interés de los alumnos en publicar y en la disponibilidad de personas revisoras (docentes y/o alumnos). Asímismo, dependiendo del interés de los alumnos podría quedar como una sola categoría o segmentarse en varias. 